%% Heyting residuation, LJ soundness, and prime-filter countermodels.

\section{Heyting algebras and IPC completeness}

\begin{annexitem}{annex:p10.heyting.1}{thm:p10-heyting-examples}{Residuation for opens and upsets}
\begin{proof}
Open subsets are closed under finite intersections and unions,
which obey the bounded distributive lattice laws.
For open \(Z,U,V\),
\[
 Z\cap U\subseteq V
 \quad\Longleftrightarrow\quad
 Z\subseteq(X\setminus U)\cup V
 \quad\Longleftrightarrow\quad
 Z\subseteq\operatorname{int}_X((X\setminus U)\cup V).
\]
The last equivalence uses that \(Z\) is open and interior is
the greatest open subset of a set.
This is precisely residuation.

Upsets are closed under unions and intersections, with bounds
\(\varnothing,W\), so their finite set operations form a
bounded distributive lattice.
The displayed implication upset is upward closed: if \(w\leq w'\),
every \(v\geq w'\) is also above \(w\), by transitivity.
If \(Z\subseteq U\Rightarrow V\) and \(w\in Z\cap U\),
reflexivity permits \(v=w\) in the defining clause, giving
\(w\in V\).
Conversely, suppose \(Z\cap U\subseteq V\).
For \(w\in Z\) and \(v\geq w\) with \(v\in U\), persistence
of \(Z\) gives \(v\in Z\), and hence \(v\in V\).
Thus \(Z\subseteq U\Rightarrow V\), proving residuation.
Neither argument uses antisymmetry.
For an empty carrier the unique subset is both bounds and
all operation values; for a nonempty carrier the bounds differ.
\end{proof}
\end{annexitem}

\begin{annexitem}{annex:p10.heyting.2}{thm:p10-heyting-completeness}{LJ soundness and the Lindenbaum Heyting algebra}
\begin{proof}
\paragraph{Algebraic soundness of the propositional rules.}
Interpret a finite LJ antecedent \(\Delta\) by the meet of
its formula values, counting the empty meet as \(1\).
The semantic assertion for a sequent is that this meet lies
below its succedent value.
Identity and bottom-left express \(a\leq a\) and \(0\leq c\).
Weakening decreases the antecedent meet; contraction preserves
it by idempotence.  Exchange is already built into the
multiset convention.
Conjunction-left is associativity of meet, and conjunction-right
uses that a common lower bound of \(a,b\) lies below \(a\wedge b\).
Each disjunction-right rule uses an inclusion into a join.
Disjunction-left follows from
\[
 (a\vee b)\wedge g=(a\wedge g)\vee(b\wedge g).
\]
Implication-right is residuation.
For implication-left, the premise inequalities \(g\leq a\)
and \(b\wedge d\leq c\) give
\[
 (a\Rightarrow b)\wedge g\wedge d
       \leq(a\Rightarrow b)\wedge a\wedge d
       \leq b\wedge d\leq c.
\]
Cut follows from \(g\leq a\) and \(a\wedge d\leq c\),
which imply \(g\wedge d\leq c\).
The displayed Mix rule has the same interpretation, since
the meet of any positive number of copies of \(a\) is \(a\).
This accounts for every rule of the propositional fragment.
Induction over a finite derivation proves soundness.

\paragraph{Congruence and quotient order.}
Reflexivity of one-premise derivability follows from identity,
and transitivity from Cut.
Its symmetric part \(\sim\) is therefore an equivalence relation.
The conjunction and disjunction rules show monotonicity of
each of these connectives in both arguments under derivability.
For implication, from \(\varphi'\vdash\varphi\) and
\(\psi\vdash\psi'\), implication-left yields
\(\varphi\to\psi,\varphi'\vdash\psi'\);
implication-right gives
\[
 \varphi\to\psi\ \vdash\ \varphi'\to\psi'.
\]
Thus implication is antitone in its first argument and
monotone in its second.
Using both directions of equivalent arguments proves
congruence for every connective.
Transitivity then makes the quotient order well-defined,
and its antisymmetry is exactly passage to \(\sim\)-classes.

The bottom-left rule makes \([\bot]\) the least element.
The formula \(\top=(\bot\to\bot)\) has an empty-premise proof,
so weakening makes \([\top]\) greatest.
Conjunction projections and conjunction-right prove the
greatest-lower-bound property of \([\varphi\land\psi]\).
Disjunction introductions and disjunction-left prove the
least-upper-bound property of \([\varphi\lor\psi]\).
Thus the quotient is a bounded lattice.

\paragraph{The residual in the quotient.}
The conjunction rules and Cut give
\[
 \chi\land\varphi\vdash\psi
       \quad\Longleftrightarrow\quad
 \chi,\varphi\vdash\psi.
\]
For example, the right-hand context derives its conjunction
by conjunction-right and weakening; cut with the left-hand
derivation gives one direction, and conjunction-left gives
the other.
Implication-right and implication-left, with Cut, give
\[
 \chi,\varphi\vdash\psi
       \quad\Longleftrightarrow\quad
 \chi\vdash\varphi\to\psi.
\]
Consequently
\[
 [\chi]\wedge[\varphi]\leq[\psi]
       \quad\Longleftrightarrow\quad
 [\chi]\leq[\varphi\to\psi].
\]
This is the required residual law for all quotient elements.

It also supplies distributivity.  In any bounded lattice
with these residuals, for arbitrary \(a,b,c,d\),
\[
\begin{aligned}
 a\wedge(b\vee c)\leq d
 &\Longleftrightarrow b\vee c\leq a\Rightarrow d\\
 &\Longleftrightarrow
       b\leq a\Rightarrow d\text{ and }c\leq a\Rightarrow d\\
 &\Longleftrightarrow
       (a\wedge b)\vee(a\wedge c)\leq d.
\end{aligned}
\]
The two elements have the same upper bounds, hence are equal.
The dual distributive identity follows from this identity
and absorption.  The quotient is therefore a Heyting algebra.

\paragraph{The canonical valuation.}
Formula induction gives
\(\llbracket\chi\rrbracket_{v_0}=[\chi]\).
An empty-premise proof of \(\chi\) can be weakened to
\(\top\vdash\chi\).
Conversely, cut \(\top\vdash\chi\) with the theorem \(\top\)
to obtain an empty-premise proof.
Thus
\[
 [\chi]=1\quad\Longleftrightarrow\quad
 \vdash_{\mathsf{IPC}}\chi.
\]
Soundness gives value \(1\) for every theorem under every
valuation in every Heyting algebra.
A non-theorem has canonical value different from \(1\),
giving completeness.
The open-set algebra of a singleton is a two-element Heyting
algebra, whose bottom value is not top.  Soundness therefore shows
that \(\bot\) is not a theorem; the quotient is nontrivial.

For finite \(\Gamma\), the conjunction rules, structural rules,
and Cut identify
\(\Gamma\vdash\varphi\) with
\((\bigwedge\Gamma)\vdash\varphi\).
Algebraic soundness gives the stated inequality in every algebra.
If the derivation fails, the canonical valuation gives
\([\bigwedge\Gamma]\not\leq[\varphi]\), refuting the inequality.
This includes the empty-premise case.
The entire quotient and soundness argument uses no choice principle
and applies to every set of propositional variables.
\end{proof}
\end{annexitem}

\begin{annexitem}{annex:p10.heyting.3}{thm:p10-heyting-prime-representation}{The prime-filter order and implication witness}
\begin{proof}
Inclusion is a partial order on \(\operatorname{Pf}(H)\).
Each \(\beta_H(a)\) is an upset.
Filter laws and primeness give preservation of \(1,0\),
binary meets, and binary joins.
If \(a\not\leq b\), the lattice filter
\(\mathord{\uparrow}a\) and ideal \(\mathord{\downarrow}b\)
are disjoint: a common member would imply \(a\leq b\).
Theorem~\ref{thm:p10-prime-filter-separation} supplies a
prime filter containing \(a\) and omitting \(b\).
Thus \(\beta_H(a)\not\subseteq\beta_H(b)\);
the map reflects order and is injective.
Taking \(a=1,b=0\) shows nonemptiness of the prime-filter
order when \(H\) is nontrivial.

\paragraph{Preserving implication.}
If \(a\Rightarrow b\in P\) and \(Q\supseteq P\) is prime
with \(a\in Q\), then
\(a\wedge(a\Rightarrow b)\in Q\).
Residuation gives \(a\wedge(a\Rightarrow b)\leq b\),
so \(b\in Q\).

For the converse, suppose \(a\Rightarrow b\notin P\).
The lattice filter generated by \(P\cup\{a\}\) is
\[
 F=\{x:\exists p\in P\quad p\wedge a\leq x\}.
\]
It contains the generators and is closed upward and under meets.
It is disjoint from the ideal \(I=\mathord{\downarrow}b\).
Indeed, a common member would give \(p\wedge a\leq b\)
for some \(p\in P\), hence \(p\leq a\Rightarrow b\)
by residuation, forcing \(a\Rightarrow b\in P\).
In particular \(F\) is proper, since \(0\in I\).
Prime-filter separation supplies \(Q\supseteq F\) disjoint
from \(I\).
This \(Q\) extends \(P\), contains \(a\), and omits \(b\).
It is precisely the witness required when implication fails.

The two directions identify \(\beta_H(a\Rightarrow b)\)
with the upset implication
\(\beta_H(a)\Rightarrow\beta_H(b)\).
Hence \(\beta_H\) is a Heyting embedding.
For trivial \(H\), no proper filter exists and
\(\operatorname{Up}(\varnothing)=\{\varnothing\}\);
the unique map is the asserted isomorphism.
For countable \(H\), every separation instance uses the
countable ZF case of the stated separation theorem.
The representation takes all prime filters as points and
defines the order by inclusion, so no simultaneous selection
of witnesses is required.
\end{proof}
\end{annexitem}

\begin{annexitem}{annex:p10.heyting.4}{thm:p10-ipc-prime-kripke-completeness}{Strong local Kripke completeness}
\begin{proof}
\paragraph{Soundness on nonempty preorders.}
For a persistent valuation on a preorder, formula induction
identifies the forcing set of each formula with its value in
the upset Heyting algebra:
the implication clause is exactly the upset residual.
All these forcing sets are upsets.
If \(\Gamma\vdash_{\mathsf{IPC}}\varphi\), a finite
supporting context has the algebraic inequality proved in
Annex~\ref{annex:p10.heyting.2}.
Every point forcing the entire premise set belongs to that
finite intersection and hence forces \(\varphi\).
This proves local soundness, using only reflexivity and
transitivity of the frame.

\paragraph{The canonical forcing lemma.}
On the partial order
\(\operatorname{Pf}(H_{\mathsf{IPC}})\), assign
\(V(p)=\beta_{H_{\mathsf{IPC}}}([p])\).
The preceding representation and formula induction give
\[
 P\forces\chi\quad\Longleftrightarrow\quad[\chi]\in P.
\]
In the implication case, preservation is exactly the
prime-filter witness theorem:
failure of \([\alpha]\Rightarrow[\beta]\) at \(P\)
has an extension containing \([\alpha]\) and omitting
\([\beta]\).

\paragraph{Separating an arbitrary premise set from its conclusion.}
Suppose \(\Gamma\not\vdash_{\mathsf{IPC}}\varphi\).
In \(H_{\mathsf{IPC}}\), let \(F_\Gamma\) be the lattice
filter generated by all \([\gamma]\), \(\gamma\in\Gamma\),
and let \(I=\mathord{\downarrow}[\varphi]\).
These are disjoint.
Otherwise a finite subset \(\Gamma_0\subseteq\Gamma\)
would satisfy
\[
 \bigwedge_{\gamma\in\Gamma_0}[\gamma]\leq[\varphi],
\]
which gives a finite-support derivation of \(\varphi\)
from \(\Gamma\) by the quotient order and the conjunction
argument above.
This contradiction includes \(\Gamma_0=\varnothing\),
whose meet is \(1\).

Prime-filter separation yields \(P\supseteq F_\Gamma\)
with \(P\cap I=\varnothing\).
The canonical forcing lemma makes every member of \(\Gamma\)
true at \(P\), and makes \(\varphi\) false there.
In particular, this countermodel has a point and its
partial-order carrier is nonempty.
It is also an allowed preorder model.
Together with soundness, this proves both stated completeness
equivalences.
An inconsistent premise set derives bottom and therefore
every formula; soundness ensures that no point forces all
its members.

\paragraph{The countable-language case in ZF.}
If \(\mathsf{Prop}\) is countable, the set of finite formulas
is countable, and so is its quotient \(H_{\mathsf{IPC}}\).
The countable version of
Theorem~\ref{thm:p10-prime-filter-separation} therefore proves
both the implication-witness and the premise-separation steps
in ZF.
The carrier is the set of all prime filters, and valuations
and order are given by explicit membership conditions.
No family of choices of prime extensions is used.
Finite or empty variable sets are included in this countable case.
\end{proof}
\end{annexitem}
